% Source: https://artofproblemsolving.com/community/c3h2579913_derangments
% Stable identity: aops-topic-2579913-q01
\textbf{(a)} Let $A_i$ be the set of seatings in which person $i$ is in their own seat. For any $k$ given people, the seatings in which all of them are in their own seats number $(n-k)!$, because the other $n-k$ people sit arbitrarily. By inclusion--exclusion,
\[\left|A_1\cup\cdots\cup A_n\right|=\sum_{k=1}^n(-1)^{k+1}\dbinom nk(n-k)!.\]
Subtracting from $n!$,
\[D_n=\sum_{k=0}^n(-1)^k\dbinom nk(n-k)!=\sum_{k=0}^n(-1)^k\dfrac{n!}{k!}. \qquad\blacksquare\]

\textbf{(b)} In a derangement, person $1$ takes some seat $j\ne1$: $n-1$ choices. Fix $j$.
\begin{itemize}
\item If person $j$ takes seat $1$, the other $n-2$ people form a derangement of their own seats: $D_{n-2}$ ways.
\item If person $j$ does not take seat $1$, then the people other than $1$ fill the seats other than $j$. Each must avoid one seat: their own seat, or seat $1$ for person $j$. Relabelling seat $1$ as $j$'s ``own'' seat, this is a derangement of $n-1$ people: $D_{n-1}$ ways.
\end{itemize}
So $D_n=(n-1)\left(D_{n-1}+D_{n-2}\right)$. Clearly $D_1=0$ and $D_2=1$. $\blacksquare$

\textbf{(c)} $D_3=2$, $D_4=9$, $D_5=44$, $D_6=265$, $D_7=1854$, and $D_8=7(1854+265)=\boxed{14833}$. This agrees with (a), and $D_n$ is the integer nearest to $\dfrac{n!}{e}$.
